\section{Lingkungan Kerja dan Struktur Program Assembly Minimal}
\label{sec:lingkungan-kerja}

Sebelum menulis program, lingkungan kerja harus disiapkan. Mata
kuliah ini menggunakan tiga perangkat lunak yang tersedia bebas,
yang dapat dipilih sesuai kebutuhan.

\subsection{Pilihan Lingkungan Kerja}

\begin{table}[H]
    \centering
    \caption{Perbandingan perangkat lunak lingkungan kerja assembly.}
    \label{tab:tools-assembly}
    \small
    \begin{tabular}{@{}p{3.5cm}p{5.2cm}p{5.1cm}@{}}
        \toprule
        \textbf{Perangkat} & \textbf{Kelebihan} &
        \textbf{Keterbatasan} \\
        \midrule
        \textbf{Emu8086} & Emulator dengan tampilan register, memori,
        dan \eng{flag} secara langsung; cocok untuk pemula &
        Menjalankan emulasi, bukan perangkat keras sungguhan \\
        \addlinespace
        \textbf{GUI Turbo Assembler} & Antarmuka grafis untuk
        TASM/TLINK dengan \eng{debugger} terpadu; mendekati alur kerja
        nyata & Perlu pemahaman struktur segmen dan linker \\
        \addlinespace
        \textbf{DOSBox} & Menyediakan lingkungan DOS asli; dapat
        menjalankan DEBUG, TASM, TLINK, dan program \file{.exe} lama &
        Antarmuka baris perintah, tanpa tampilan register otomatis \\
        \bottomrule
    \end{tabular}
\end{table}

\begin{aktivitasbox}
Langkah penyiapan yang harus diselesaikan mahasiswa sebelum
pertemuan berikutnya:
\begin{enumerate}
    \item Pasang Emu8086, lalu buat satu berkas program sederhana dan
          jalankan.
    \item Pasang GUI Turbo Assembler beserta TASM dan TLINK, lalu
          rakit ulang program yang sama.
    \item Pasang DOSBox, tambahkan perintah \file{mount} ke direktori
          kerja, dan jalankan perintah \file{DEBUG} di dalamnya.
    \item Dokumentasikan langkah pemasangan dalam jurnal praktikum.
\end{enumerate}
\end{aktivitasbox}

\subsection{Struktur Program Assembly Minimal}

Program assembly 8086 untuk lingkungan DOS memiliki kerangka baku.
Kerangka tersebut tampak pada Kode~\ref{lst:kerangka-program}.

\begin{lstlisting}[
    language=MASM,
    caption={Kerangka baku program assembly 8086 untuk DOS.},
    label={lst:kerangka-program}
]
.MODEL SMALL        ; tentukan model memori
.STACK 100h         ; sediakan ruang stack

.DATA               ; ------- segmen data -------
    ; deklarasi variabel dan konstanta di sini

.CODE               ; ------- segmen kode -------
START:              ; label titik masuk program
    MOV AX, @DATA   ; inisialisasi register DS
    MOV DS, AX

    ; instruksi program di sini

    MOV AH, 4Ch     ; keluar dari program
    MOV AL, 00h
    INT 21h

END START           ; akhiri dan tunjuk titik masuk
\end{lstlisting}

Setiap bagian kerangka tersebut memiliki peran sebagai berikut.

\begin{description}
    \item[\dir{.MODEL SMALL}] Menetapkan model memori. Model
          \file{SMALL} berarti program memakai satu segmen kode dan
          satu segmen data, masing-masing berukuran maksimum 64 KB.
          Model ini paling banyak dipakai untuk program pembelajaran.

    \item[\dir{.STACK 100h}] Menyediakan ruang \eng{stack} sebesar
          \hex{100} byte (\dec{256} byte) untuk menyimpan alamat
          kembali pemanggilan prosedur dan data sementara.

    \item[\dir{.DATA}] Awal segmen data. Semua variabel dan konstanta
          dideklarasikan di sini menggunakan direktif \dir{DB},
          \dir{DW}, atau \dir{DD}.

    \item[\dir{.CODE}] Awal segmen kode. Seluruh instruksi yang
          dieksekusi berada pada bagian ini.

    \item[\file{START:}] Label yang menandai instruksi pertama. Nama
          label bebas, tetapi harus sama dengan nama yang ditulis
          setelah \dir{END}.

    \item[\dir{END START}] Menutup berkas program sekaligus
          memberi tahu assembler bahwa eksekusi dimulai dari label
          \file{START}.
\end{description}

\subsection{Inisialisasi Register Segmen}

Program \file{.EXE} tidak otomatis memiliki \reg{DS} yang benar. Dua
instruksi berikut wajib ada di awal program:

\begin{lstlisting}[language=MASM, numbers=none]
MOV AX, @DATA       ; salin alamat segmen data ke AX
MOV DS, AX          ; lalu pindahkan ke DS
\end{lstlisting}

Mengapa harus melalui \reg{AX} dan tidak langsung
\lstasm|MOV DS, @DATA|? Karena arsitektur 8086 melarang pemindahan
nilai segera (\eng{immediate}) langsung ke register segmen.
Pemindahan ke register segmen hanya boleh berasal dari register
umum. Aturan ini akan dibahas lebih rinci pada Bab~\ref{chap:register}
dan Bab~\ref{chap:translasi}.

\begin{infobox}
Simbol \code{@DATA} adalah simbol bawaan yang dihasilkan oleh
\dir{.MODEL}. Simbol ini berisi nomor segmen tempat \dir{.DATA}
berada. Bila program tidak memakai \dir{.MODEL}, nama segmen harus
ditulis manual melalui direktif \dir{SEGMENT}.
\end{infobox}

\subsection{Alur Eksekusi Program Pertama}

Ketika program dijalankan, urutan kejadiannya adalah sebagai berikut.

\begin{enumerate}
    \item DOS memuat berkas \file{.exe}, membuat \eng{Program Segment
          Prefix}, dan menetapkan \reg{CS} serta \reg{IP} ke titik
          masuk program.
    \item Instruksi pertama menginisialisasi \reg{DS} agar penunjukan
          variabel pada segmen data menjadi benar.
    \item Instruksi program dijalankan satu per satu sesuai urutan.
    \item Layanan DOS \code{INT 21h} fungsi \hex{09} mencetak string
          ke layar.
    \item Layanan DOS \code{INT 21h} fungsi \hex{4C} mengakhiri
          program dan mengembalikan kendali ke DOS.
\end{enumerate}
